\setcounter{section}{-1}
\section{\S 3.0 氧化还原反应复习}\label{sect:氧化还原反应复习}

    \subsection{氧化还原反应的基本概念}\label{subsec:氧化还原基本概念}

        \textcolor{red}{特征}：反应前后，元素化合价有升降。

        \textcolor{red}{本质}：电子的转移（得失或偏移）。

        有电子转移的化学反应就是氧化还原反应。

    \subsection{氧化剂与还原剂}\label{subsec:氧化剂与还原剂}

        \begin{center}
        \begin{tabular}{|c|c|c|c|}
            \hline
            \textbf{概念} & \textbf{特征} & \textbf{反应类型} & \textbf{产物} \\
            \hline
            氧化剂 & 得电子（降价） & 还原反应 & 还原产物 \\
            \hline
            还原剂 & 失电子（升价） & 氧化反应 & 氧化产物 \\
            \hline
        \end{tabular}
        \end{center}

        \textcolor{red}{口诀：得（电子）降（价）还（原反应）——氧化剂；失（电子）升（价）氧（化反应）——还原剂。}

    \subsection{氧化还原基本规律}\label{subsec:氧化还原基本规律}

        \subsubsection{守恒律}\label{subsubsec:守恒律}

            氧化剂得电子总数 \uline{=} 还原剂失电子总数（化合价升降守恒）。

        \subsubsection{强弱律}\label{subsubsec:强弱律}

            \[\cemh{\text{氧化剂} + \text{还原剂} \xleq{} \text{还原产物} + \text{氧化产物}}\]

            氧化性：氧化剂 \textcolor{red}{>} 氧化产物

            还原性：还原剂 \textcolor{red}{>} 还原产物

        \subsubsection{转化律}\label{subsubsec:转化律}

            同种元素高低价态之间的反应，氧化产物和还原产物可能为同一（或不同一）中间价态。但升价元素反应后的化合价不高于降低元素反应后的化合价。

            \textcolor{red}{只靠拢，不交叉。}

        \subsubsection{先后律}\label{subsubsec:先后律}

            相同条件下，还原性更强或氧化性更强的物质先发生反应。

    \subsection{影响氧化还原反应的因素}\label{subsec:影响氧化还原反应的因素}

        \textcolor{red}{物质的氧化性、还原性会受浓度、pH、温度的影响。}

        \begin{xmp}[FeI 氧化还原]{浓度对氧化还原方向的影响}
            \cemh{2Fe^{3+} + 2I^- <=> 2Fe^{2+} + I2}

            \textcolor{red}{\(K = \num{6.2E7}\)}

            \cemh{Fe^{3+}} 浓度越大，\cemh{Fe^{3+}} 氧化性就越强；\cemh{I2} 同理。

            物质的氧化性（还原性）强弱，是\textcolor{red}{等浓度等条件}下的比较。如果氧化性（还原性）差异不大，浓度可能会改变二者的氧化性或还原性强弱关系。
        \end{xmp}

        \begin{xmp}[H 氧化性]{\cemh{H+} 对氧化性的影响}
            半反应：\cemh{MnO2 + 2e- + 4H+ \xleq{} Mn^{2+} + 2H2O}

            \cemh{H+} 的浓度越大，\cemh{MnO2} 的氧化性越强。

            实验室制 \cemh{Cl2}：加热使 \cemh{Cl2} 逸出，\cemh{Cl-} 的还原性增强，同时 \cemh{Cl2} 浓度减小氧化性减弱。
        \end{xmp}

        \begin{xmp}[王水]{王水溶金}
            王水（浓盐酸和浓硝酸）溶金时会有 \cemh{NO} 气体逸出，溶液中金元素主要以 \cemh{[AuCl4]-} 形式存在。

            浓盐酸的作用：\textcolor{red}{提供 \cemh{Cl-}，降低 \cemh{Au} 的还原性（形成配合物）；增大 \cemh{Cl-} 浓度。}

            把浓盐酸换成 \cemh{NaCl} 溶液后，还能溶金吗？\textcolor{red}{不能，需要 \cemh{H+} 的氧化性环境。}
        \end{xmp}

    \subsection{常见氧化剂与还原剂}\label{subsec:常见氧化剂与还原剂}

        \subsubsection{常见氧化剂}\label{subsubsec:常见氧化剂}

            \begin{itemize}
                \item 活泼非金属单质：\cemh{F2}、\cemh{Cl2}、\cemh{Br2}、\cemh{O2}、\cemh{O3}
                \item 高价金属阳离子：\cemh{Fe^{3+}}、\cemh{Cu^{2+}}
                \item 高氧化态含氧酸或含氧酸根：\cemh{KMnO4}、\cemh{K2Cr2O7}、\cemh{HNO3}、\cemh{HClO}、\cemh{NaClO}、\cemh{H2O2}
                \item 其他：\cemh{MnO2}、\cemh{Na2O2}
            \end{itemize}

        \subsubsection{常见还原剂}\label{subsubsec:常见还原剂}

            \begin{itemize}
                \item 活泼金属单质：\cemh{Na}、\cemh{Mg}、\cemh{Al}、\cemh{Zn}、\cemh{Fe}
                \item 某些非金属单质：\cemh{H2}、\cemh{C}、\cemh{Si}
                \item 低价金属阳离子：\cemh{Fe^{2+}}
                \item 低价阴离子或化合物：\cemh{I-}、\cemh{S^{2-}}、\cemh{SO3^{2-}}、\cemh{H2S}、\cemh{CO}
            \end{itemize}

    \subsection{半反应的书写}\label{subsec:半反应书写}

        所有的氧化还原反应都可以拆成两个半反应：

        \begin{itemize}
            \item \textcolor{red}{氧化反应}（失电子）：还原剂 \makebox[6em]{\cemh{- e- \xleq{}}} 氧化产物
            \item \textcolor{red}{还原反应}（得电子）：氧化剂 \makebox[6em]{\cemh{+ e- \xleq{}}} 还原产物
        \end{itemize}

        \subsubsection{书写步骤}\label{subsubsec:半反应书写步骤}

            \begin{enumerate}
                \item 确定电极反应物和产物
                \item 确定得失电子数目（根据化合价升降）
                \item 根据溶液酸碱性补 \cemh{H+} 或 \cemh{OH-}（或 \cemh{H2O}）
                \item 检查原子守恒、电荷守恒
            \end{enumerate}

            \begin{center}
            \small
            \begin{tabular}{|c|c|c|}
                \hline
                \textbf{条件} & \textbf{左边多氧} & \textbf{左边少氧} \\
                \hline
                酸性介质 & \cemh{+ 2H+ \xleq{} H2O} & \cemh{+ H2O \xleq{} 2H+} \\
                \hline
                碱性介质 & \cemh{+ H2O \xleq{} 2OH-} & \cemh{+ 2OH- \xleq{} H2O} \\
                \hline
            \end{tabular}
            \end{center}

        \begin{xmp}[半反应书写例 1]{从总反应拆分}
            将反应 \cemh{2FeCl3 + Cu \xleq{} CuCl2 + 2FeCl2} 拆写为半反应式。

            氧化反应：\cemh{Cu - 2e- \xleq{} Cu^{2+}}

            还原反应：\cemh{Fe^{3+} + e- \xleq{} Fe^{2+}}
        \end{xmp}

        \begin{xmp}[半反应书写例 2]{硫化氢与 \cemh{FeCl3} 反应}
            将硫化氢气体通入 \cemh{FeCl3} 溶液所发生的反应拆写为半反应式。

            氧化反应：\cemh{H2S - 2e- \xleq{} S v + 2H+}

            还原反应：\cemh{Fe^{3+} + e- \xleq{} Fe^{2+}}
        \end{xmp}

        \begin{xmp}[半反应书写例 3]{铜与稀硝酸}
            将 \cemh{Cu} 与稀硝酸的反应拆写为半反应式。

            氧化反应：\cemh{Cu - 2e- \xleq{} Cu^{2+}}

            还原反应：\cemh{NO3- + 3e- + 4H+ \xleq{} NO ^ + 2H2O}
        \end{xmp}

        \begin{xmp}[半反应书写例 4]{\cemh{SO2} 与 \cemh{FeCl3}}
            \cemh{FeCl3} 溶液中通入 \cemh{SO2} 气体后溶液颜色变浅，滴入 \cemh{BaCl2} 溶液有白色沉淀产生。

            氧化反应：\cemh{SO2 + 2H2O - 2e- \xleq{} SO4^{2-} + 4H+}

            还原反应：\cemh{Fe^{3+} + e- \xleq{} Fe^{2+}}

            总反应：\cemh{2Fe^{3+} + SO2 + 2H2O \xleq{} 2Fe^{2+} + SO4^{2-} + 4H+}
        \end{xmp}

        \begin{xmp}[半反应书写例 5]{碱性条件}
            已知在 \cemh{KOH} 溶液中发生反应：
            \cemh{Ag2O + Zn + H2O \xleq{} 2Ag + Zn(OH)2}
            写出该反应的半反应式。

            氧化反应：\cemh{Zn - 2e- + 2OH- \xleq{} Zn(OH)2}

            还原反应：\cemh{Ag2O + 2e- + H2O \xleq{} 2Ag + 2OH-}
        \end{xmp}
